\section{Related Work} \label{ch5:sec:Related}
Energy–performance optimization in HPC systems has been studied extensively across scheduling theory, dynamic voltage and frequency scaling (DVFS), power capping, and learning-based control. Prior work can broadly be categorized into offline scheduling and speed-scaling theory, control-theoretic and model-based runtime power management, learning-based power control, and application-aware power modeling and tuning. Our work differs from these approaches by introducing a preference-driven, job-level controller that operates independently of the scheduler and avoids retraining for different energy–performance trade-offs.

Classical work on energy-aware scheduling formulates the problem as executing a set of jobs on a variable-speed processor to minimize energy consumption. In the strict-deadline setting, each job is defined by a release time $r$, deadline $d$, and workload $w$, representing the required number of CPU cycles. The optimal offline solution is given by the Yao–Demers–Shenker (YDS) algorithm, which identifies intervals of maximum workload density and assigns constant execution speeds within those intervals \cite{YDS}. Earliest Deadline First (EDF) is commonly used as an online heuristic in this setting \cite{YDS}.

For scenarios without explicit deadlines, prior work considers minimizing a combination of energy and flow time, leading to different optimality criteria and competitive guarantees \cite{ALBERS20151194}. Online algorithms are often evaluated using competitive analysis, where an algorithm is said to be c-competitive if its objective value is within a factor 
$c$ of the optimal offline solution for the same input \cite{YDS,ALBERS20151194}. While these formulations provide strong theoretical insights, they assume precise knowledge of job characteristics and speed–power relationships, and they do not directly translate to modern HPC systems where applications exhibit complex, hardware-dependent behavior.

DVFS has long been a primary mechanism for reducing energy consumption while meeting performance requirements. Control-theoretic approaches model the relationship between frequency, performance, and power, and use feedback control to regulate execution speed \cite{DVFS_Control, DVFS_optimization}. Several works demonstrate that DVFS-based controllers can effectively meet timing constraints with reduced energy, particularly for compute-bound workloads \cite{DVFS_Control}.

However, recent processor designs increasingly move DVFS control into hardware, limiting the ability of software to directly set operating frequencies. This trend has motivated the use of hardware power capping interfaces such as Intel RAPL, where software specifies a power limit and the hardware manages voltage and frequency internally \cite{CoPPer}. Control-based approaches using power caps have been proposed to meet soft real-time performance targets, but they must deal with non-linear and diminishing returns in the power–performance relationship \cite{CoPPer}. These methods typically rely on adaptive control mechanisms and runtime feedback, and are often designed for single objectives or fixed performance targets.

Machine learning and reinforcement learning (RL) have been explored to address the limitations of fixed models and hand-tuned controllers. Early work applies Q-learning and Double-Q learning to DVFS governors, enabling systems to adapt frequency decisions based on observed workload behavior \cite{DDQL}. These approaches demonstrate improved energy savings compared to static governors, but they require online exploration and are sensitive to training stability and convergence time.

More recent efforts employ model-based or multi-objective learning techniques to predict performance and energy under different power cap configurations \cite{MOML}. These methods often construct performance and energy models using micro-benchmarks or application features and then identify Pareto-optimal configurations. While effective, such approaches typically require extensive profiling, retraining for new architectures or applications, and separate optimization runs for different objectives or user preferences.

Offline RL has been proposed as a way to avoid unsafe or costly online exploration by training policies from pre-collected data \cite{offlineRL}. However, most existing RL-based approaches still produce a single policy optimized for a fixed objective, requiring multiple agents or retraining when the desired trade-off changes.

To make power management practical, many studies rely on HPC proxy applications and benchmark suites to capture representative computation and communication patterns\cite{MOML,HPC_proxy_applications}. Proxy applications enable controlled experiments to study how performance degrades under power caps and how energy efficiency varies across workloads. Empirical studies show that different applications respond very differently to power limiting, highlighting the need for application-aware control \cite{MOML}.

Several works combine proxy-based profiling with analytical or data-driven models to guide power allocation decisions at the node or cluster level \cite{MOML,OS_noise}. While these approaches provide valuable insights into workload behavior, they often assume static optimization goals and do not expose an intuitive mechanism for users to express acceptable performance degradation.

In contrast to prior work, our approach focuses on job-level runtime control rather than scheduler-level resource allocation or cluster-level power distribution. We propose a preference-driven multi-objective reinforcement learning framework that allows users to specify their desired trade-off between performance and energy using a simple preference parameter. By leveraging mathematical models of application behavior to generate training data, we eliminate the need for interpolation models and enable efficient offline training.

Unlike existing learning-based approaches, our method produces a single controller that can operate across a range of user-defined performance degradation levels, avoiding the need to train separate agents for different objectives. This design complements existing schedulers by providing fine-grained, application-aware power control during job execution, while preserving scheduler autonomy over resource allocation and job ordering.